\section{What Comes Out First}
\label{sec:ch08-what-comes-out-first}
\index{extraction!choosing the first domain|(}%

Mosaic has six domains and one of them can leave without asking permission.
Catalog stores a product's identity, its stock keeping unit, its title, its
description and its category, and it reads nothing else. Pricing reads a product
identifier. Inventory reads a product identifier. Reviews reads a product
identifier and a customer identifier. Ordering reads all three. Point the arrows
and one folder has none coming out of it.

\index{extraction!acyclic dependencies}%
That is the rule I would apply to any decomposition, and it is duller than the
architecture diagrams suggest: extract the domain that references nobody, and
the extraction is a move. Extract one in the middle and it is a negotiation
about who holds what. Chapter~\ref{ch:strangling-the-monolith} carves the
remaining five, where the arrows do point outward and the answer is
\code{@override} and a migration rather than a rename.

Figure~\ref{fig:ch08-the-cut} is the shape of the change.

\begin{figure}[htbp]
  \centering
  % One service becomes two, and Product is declared on both sides of the
% boundary with the same key. Referenced from chapter 08. Bare tikzpicture; the
% figure environment, caption and label live at the call site.
\begin{tikzpicture}[
    font=\small,
    svc/.style={draw, rounded corners=2pt, minimum width=46mm,
                minimum height=9mm, align=center},
    fld/.style={draw, rounded corners=2pt, minimum width=14mm,
                minimum height=5mm, font=\scriptsize, align=center},
    typ/.style={draw, rounded corners=2pt, minimum width=46mm,
                minimum height=19mm, align=left, font=\scriptsize},
    boundary/.style={gray, dashed, thick},
    shared/.style={{Stealth[length=2mm]}-{Stealth[length=2mm]}, gray},
    feed/.style={-{Stealth[length=2mm]}},
    lbl/.style={font=\scriptsize, align=center, inner sep=1pt}
  ]

  % the process boundary
  \draw[boundary] (6.35,0.9) -- (6.35,-5.4);
  \node[lbl, gray, anchor=south] at (6.35,1.0) {two processes};

  % -- Catalog, on the left ------------------------------------------------
  \node[svc] (cat) at (3.05,0.2)
    {\code{Mosaic.Catalog}\\\scriptsize :5101, database \code{catalog}};

  \node[fld] (catfld) at (3.05,-1.85) {\code{Catalog/}};

  \node[typ] (catp) at (3.05,-4.0)
    {\code{type Product @key(fields: "id")}\\[2pt]
     \code{id}\\
     \code{sku}\\
     \code{title}\\
     \code{description}\\
     \code{category}};

  \draw[feed] (catfld) -- (3.05,-3.05);

  % -- Mosaic, on the right ------------------------------------------------
  \node[svc] (mos) at (9.65,0.2)
    {\code{Mosaic.Api}\\\scriptsize :5100, database \code{mosaic}};

  \node[fld] (pri) at (9.65,-1.2) {\code{Pricing/}};
  \node[fld] (inv) at (9.65,-1.85) {\code{Inventory/}};
  \node[fld] (rev) at (9.65,-2.5) {\code{Reviews/}};

  \node[typ] (mosp) at (9.65,-4.0)
    {\code{type Product @key(fields: "id")}\\[2pt]
     \code{id}\\
     \code{price}\\
     \code{availableQuantity}\\
     \code{reviews}\\
     \code{averageRating}};

  \draw[feed] (rev) -- (9.65,-3.05);

  % -- the one thing they share --------------------------------------------
  \draw[shared] (5.35,-4.0) -- (7.35,-4.0);
  \node[lbl, anchor=north, align=center] at (6.35,-5.5)
    {one type, because both sides\\spell the name and the key the same way};

\end{tikzpicture}

  \caption{One deployable becomes two. \code{Product} is declared on both sides
  of the line with the same key, and four of its fields answer from a process
  that has never seen a product row. Ordering is left out of the picture
  because it contributes nothing to \code{Product}; it only holds identifiers,
  which is section~\ref{sec:ch08-extending-a-foreign-entity}. The dashed boundary is
  the only new thing here. Every folder was already where it is.}
  \label{fig:ch08-the-cut}
\end{figure}

\subsection*{What git called it}

\index{extraction!as a file move}%
Eight files changed path, and git will tell you how much of each survived the
trip:

\begin{minted}[fontsize=\footnotesize,breaklines]{text}
rename src/{Mosaic.Api => Mosaic.Catalog}/Catalog/CatalogRegistration.cs (81%)
rename src/{Mosaic.Api => Mosaic.Catalog}/Catalog/Data/CatalogDataLoaders.cs (87%)
rename src/{Mosaic.Api => Mosaic.Catalog}/Catalog/Data/CatalogSeedData.cs (98%)
rename src/{Mosaic.Api => Mosaic.Catalog}/Catalog/Data/CatalogService.cs (68%)
rename src/{Mosaic.Api => Mosaic.Catalog}/Catalog/Data/ProductConfiguration.cs (96%)
rename src/{Mosaic.Api => Mosaic.Catalog}/Catalog/Model/ProductCategory.cs (73%)
rename src/{Mosaic.Api => Mosaic.Catalog}/Catalog/Types/CatalogQueries.cs (97%)
rename src/{Mosaic.Api => Mosaic.Catalog}/Catalog/Types/ProductNode.cs (55%)
\end{minted}

Five of the eight changed nothing but their \code{namespace} line and their
\code{using} lines: two lines each for the registration, the DataLoader, the
seed data and the entity configuration, one for the enum.
\code{CatalogQueries} changed three, and not one of them was a root field. The
four that chapters~\ref{ch:hotchocolate-quickly}
to~\ref{ch:schema-design-that-survives-change} built, including the deprecation
chapter~\ref{ch:schema-design-that-survives-change} put on \code{products},
compiled in the new project untouched.

The percentages are worth reading rather than skipping. Ninety-eight per cent
for a file of seed data whose only change was where it says it lives.
Fifty-five for \code{ProductNode}, which kept both of its methods and gained
paragraphs of comment about what is deliberately \emph{not} in it, which
section~\ref{sec:ch08-two-attributes} is about.

\code{CatalogService} is the sixty-eight, and it lost one thing:

\begin{minted}{csharp}
public sealed class CatalogService(CatalogDbContext db)
\end{minted}

\index{ServiceCallCounter@\code{ServiceCallCounter}!left behind}%
It used to take two things: the shared \code{MosaicDbContext}, and a
\code{ServiceCallCounter}. The context is Catalog's own now, and the counter is
gone along with the calls to it, which is why all five of its query methods
collapsed into expressions with no body left to hold. That counter was built in
chapter~\ref{ch:the-life-of-a-request} to count what one process was doing to
itself, and it belongs to Mosaic. A service that answers one domain's questions
has nothing to count with it, and I would rather a new service start with no
instrumentation than with somebody else's.

\index{extraction!the type that did not move}%
One file in that folder did not move at all, and its absence from the list is
the whole of section~\ref{sec:ch08-extending-a-foreign-entity}: \code{Model/Product.cs}
exists on both sides now, as two different types.

\subsection*{The three files that did not change}

\index{ObjectType@\code{[ObjectType<T>]}!across a service boundary}%
\code{Product.price} is a resolver in \code{Pricing/Types/}.
\code{Product.availableQuantity} is one in \code{Inventory/Types/}.
\code{Product.reviews} and \code{Product.averageRating} are two more in
\code{Reviews/Types/}. All four hang off a type declared in another folder,
through an \code{[ObjectType<Product>]} partial class, because
chapter~\ref{ch:hotchocolate-quickly} made that the rule for this repository and
held to it even where a shortcut was tempting.

None of those three files changed by a character when Catalog left. The type
they extend is not in the same assembly any more. They do not know that.

\index{decomposition!what the layout rule buys}%
I want to be precise about what that buys, because it is easy to oversell. It
does not make the extraction free: there is a new project, a new database, a new
schema to publish, and section~\ref{sec:ch08-only-breaks-together}'s three
surprises to work through. What it buys is that the resolvers stayed put. A
version of Mosaic that had registered every field on \code{Product} in one
central place would have needed that list unpicked, and every field on the wrong
side of the cut rewritten against a type it could no longer see. That is the
rewrite the layout rule was insurance against, and this is the chapter that
collects.

\subsection*{Two databases, and nothing in the compose file}

\index{extraction!splitting the database}%
Catalog took its table with it. \code{CatalogDbContext} is
\code{MosaicDbContext} with five of its six \code{DbSet} properties deleted, and
\code{ProductConfiguration} moved into the new assembly, where
\code{ApplyConfigurationsFromAssembly} finds it exactly as before.

The new service points at a database called \code{catalog} on the same
PostgreSQL container. That needed no change to \code{docker-compose.yml} and no
initialisation script, because Entity Framework Core's
\code{EnsureCreatedAsync} creates the database as well as the schema and the
container's user is allowed to. Two services on one server is a deployment
detail; two services on one table would be a design failure, and this is the
first of those rather than the second.

Nothing in the seed data had to move either. Every domain already carried its
own copy of the identifier pattern:

\begin{minted}{csharp}
private static Guid ProductId(int n) => new($"a0000000-0000-4000-8000-{n:D12}");
\end{minted}

That line appears in Pricing's seed data, Inventory's, Reviews' and Ordering's,
each with a comment saying the duplication is on purpose. In
chapter~\ref{ch:hotchocolate-quickly} it read as a slightly fussy way to avoid
sharing a file. Here it is the reason a price still finds its product across a
process boundary.

The code, then, was the easy half. The schema is where the chapter lives.

\index{extraction!choosing the first domain|)}%
